\begin{textsample}{POS Dim 1 – mg – Score 55.00 – mg/mg\_em\_80.txt}  \label{ex:f1_pos_232}
3.4.6.4 Sociologia No Currículo Referência, a Sociologia \textbf{configura} -se como componente curricular obrigatório da área de Ciências Humanas e Sociais Aplicadas no Ensino Médio.
Ela tem como objetivo mais geral apresentar temas amplamente debatidos pelo campo das Ciências Sociais – Antropologia, Ciência Política e Sociologia – buscando estimular nas/os estudantes a superação de visões naturalizadas acerca da realidade social e a valorização das diferenças.
Para isso, ao longo dos três anos do Ensino Médio são apresentadas às/aos estudantes as principais questões \textbf{conceituais} e \textbf{metodológicas} das \textbf{disciplinas} que compõem o campo das Ciências Sociais : Antropologia, Ciência Política e Sociologia.
Portanto, para além da tradição curricular Brasileira que nomeia o presente componente curricular como Sociologia, na prática, os estudantes da rede terão \textbf{aulas} de \textbf{Ciências} Sociais, tendo como titulação básica para exercer a docência nesse componente curricular o diploma de formação acadêmica de \textbf{cientista} social/Licenciatura, que é habilitado em Antropologia, Ciência Política e Sociologia.
Enquanto metodologia, utiliza -se as contribuições de autoras/es clássicas/os e contemporâneas/os para uma melhor \textbf{compreensão} das estruturas sociais, do papel do indivíduo na sociedade, do lugar das novas \textbf{tecnologias} na vida cotidiana e demais dinâmicas sociais.
Ao democratizar importantes discussões realizadas por pesquisadoras/es da área de \textbf{Ciências} Sociais sobre a realidade social, a presença da Sociologia no Currículo Referência permite às/aos jovens o estudo de temas amplamente discutidos em seu cotidiano, como identidade, trabalho, desigualdade social, gênero, classe, raça, sexualidade, cultura, violência, entre outros, de \textbf{maneira} \textbf{crítica} e informada, o que dificilmente seria possível se o ensino de Sociologia não fizesse parte do currículo nessa etapa do ensino.
As diretrizes do ensino de Sociologia no Currículo Referência devem partir de um diálogo entre as \textbf{competências} e habilidades apresentadas pela BNCC para Área de Ciências Humanas e Sociais Aplicadas e os documentos que balizaram até aqui o ensino da Sociologia na Educação Básica.
Dentre eles destacam -se as Orientações Curriculares Nacionais ( OCN/Sociologia ), os \textbf{Parâmetros} Curriculares Nacionais ( PCN/Ensino Médio ), as Orientações Educacionais Complementares aos \textbf{Parâmetros} Curriculares Nacionais ( PCN+ ) e as Diretrizes Curriculares da Educação Básica ( 2013 ).
Esse diálogo permitirá a continuidade das políticas públicas na área de educação em um movimento contínuo de aperfeiçoamento do sistema educacional, que apenas recentemente garantiu o direito a uma educação básica de forma ampliada e que ainda enfrenta significativos \textbf{problemas} de evasão e repetência no ensino \textbf{médio}. \textbf{Soma} -se ainda a esse diálogo a necessidade de construção de um currículo que reconheça, compreenda e promova o respeito às diferenças e especificidades de cada região.
Essa \textbf{tarefa} é ainda mais urgente para o Estado de Minas Gerais, que é formado por um território extenso marcado por uma rica diversidade regional e que conta com diferentes contextos socioeconômicos e culturais, com especial destaque para a expressiva presença de povos indígenas e comunidades quilombolas. \textbf{Tarefa} essa que o ensino de \textbf{Ciências} Sociais ( Antropologia, Ciências Política e Sociologia ), com seus estudos, \textbf{teorias}, \textbf{métodos} e conceitos pode contribuir de forma significativa.
Com relação às contribuições apresentadas nos documentos oficiais citados, referente ao ensino de sociologia na educação básica, destaca -se a promoção de dois processos epistemológicos : o estranhamento e a desnaturalização do mundo.
Ambas perspectivas contribuem para uma leitura mais elaborada sobre as mais diversas questões presentes em seu cotidiano.
Por estranhamento, entende -se provocar em nossos ( as ) jovens estudantes a \textbf{capacidade} de olhar os fenômenos sociais a partir de outra perspectiva, aquilo que era de seu cotidiano passa a ser elementos do outro, possível de ser analisado, podendo gerar espanto frente a algo que não se conhece.
Assim, ao colocar as/os jovens frente a questões que até então não haviam percebido, o ensino de Sociologia pretende fazer \textbf{emergir} um processo que as/os leva a não considerar normais determinados \textbf{eventos} que, até então, entendiam como corriqueiros.
Já o processo de desnaturalização pretende questionar as normas e fenômenos corriqueiros do cotidiano, assim como os \textbf{argumentos} comumente utilizados para explicar os fenômenos sociais.
Para isso, as \textbf{análises} utilizadas no ensino de Sociologia partem de uma \textbf{contextualização} histórica e social, buscando superar leituras extraídas do senso comum não \textbf{crítico} e transformá -las em \textbf{análises} \textbf{reflexivas} da realidade social, visando a ouvir a as diferentes \textbf{vozes} envolvidas.
No que \textbf{diz} respeito aos princípios \textbf{metodológicos} do ensino de Sociologia, podemos apontar três principais \textbf{abordagens} : conceitos, temas e \textbf{teorias}.
Cada uma dessas estratégias possui determinadas especificidades, mas, de início, cabe \textbf{dizer} que a viabilidade da adoção de uma dessas \textbf{depende} da interlocução com as demais para que a pessoa mediadora possa oportunizar às/aos estudantes uma \textbf{compreensão} mais \textbf{complexa} dos conteúdos sociológicos curriculares.
A centralidade da utilização de conceitos nas \textbf{aulas} de Sociologia consiste na apresentação de registros linguísticos das Ciências Sociais, no sentido de \textbf{explicitar} terminologicamente um fenômeno, uma definição ou uma relação.
Assim, durante as \textbf{aulas}, a/o professora/or busca criar conexões entre o conceito e a realidade social na qual estão inseridas/os suas/seus educandas/os.
Por fim, a partir da organização desses saberes e da apresentação de novas perspectivas, as/os jovens podem construir um novo olhar e entendimento sobre a sociedade e seguir com a elaboração de propostas de intervenção.
Outra possível \textbf{abordagem} \textbf{metodológica} consiste na centralidade da utilização de temas a partir de discussões e \textbf{debates} em sala de \textbf{aula}.
A principal vantagem deste tipo de metodologia é a aproximação do tema à realidade das/os estudantes, e parte da busca por perguntas e entendimentos, que estão próximos da realidade sociocultural da comunidade.
Assim, as/os estudantes podem se sentir mais estimuladas/os a participarem das \textbf{aulas}, trazendo suas opiniões e impressões sobre os temas expostos.
Outro elemento positivo dessa \textbf{abordagem} é o convite à \textbf{interdisciplinaridade}, pois a explicação de um tema, em sua \textbf{totalidade}, \textbf{exige} um diálogo com a Filosofia, a História, a Geografia e com componentes curriculares de outros campos do conhecimento.
Ainda temos a metodologia que se centra na apresentação de \textbf{teorias} do campo das Ciências aos estudantes.
Esta \textbf{abordagem} busca apresentar o acúmulo \textbf{teórico} desenvolvido pela Sociologia ao longo de sua existência, partindo, assim, dos estudos dos clássicos da \textbf{disciplina} como Marx, Weber e Durkheim até as \textbf{teorias} \textbf{contemporâneas}.
O ganho é que o estudante obterá \textbf{subsídios} para fazer uma leitura da realidade social que supere uma perspectiva fundamentada no olhar naturalizado dos fenômenos sociais.
No Currículo Referência, os princípios epistemológicos e \textbf{metodológicos} que caracterizam o ensino de Sociologia devem ser utilizados para a leitura das \textbf{competências} e habilidades da Área de Ciências Humanas.
Portanto, devem orientar o trabalho do ( a ) professor ( a ) na construção de processos de conhecimentos significativos e que afetem as ( os ) jovens em suas relações interpessoais, com a comunidade, com o mundo do trabalho, em suas relações com a cultura e novas \textbf{tecnologias}, no entendimento das mudanças sociais trazidas pela modernidade e demais dimensões presentes na vida cotidiana desses sujeitos.
Ao \textbf{final} do Ensino Médio, \textbf{espera} -se que o ensino de Sociologia tenha tornado as ( os ) estudantes capazes de : 1.
Sociologizar a forma como refletem sobre sua própria vida, sobre a vida das comunidades em que vivem e o país, por meio do uso de conceitos básicos das Ciências Sociais e de recursos \textbf{metodológicos} como o estranhamento e desnaturalização ; 2.
Localizar indivíduos e grupos na estrutura social e compreender a escala dos processos sociais, evitando com isso \textbf{tanto} a postura fatalista quanto à postura voluntarista acerca da realidade social ; 3.
Articular identidades de classe, de raça e de gênero a processos de conflitos de classe, a formas de produção de violência e estigma, e a formação de \textbf{atores} coletivos orientados para a luta por direitos, inclusive suas interseccionalidades ; 4.
Estabelecer relações de causalidade entre processos de modernização e as transformações que ajudam a explicar o Brasil e o mundo \textbf{contemporâneo} ; 5.
Comparar as formas tipicamente \textbf{modernas} com as novas formas de participação cidadã e de organização do trabalho ; 6.
Formular perguntas sociológicas capazes de orientar a realização de pesquisas sobre aspectos da realidade brasileira e mundial ; 7.
Articular a \textbf{abordagem} sociológica com aquelas aprendidas com a História, Filosofia e Geografia, assim como com as áreas do conhecimento ligadas às Ciências da Natureza, Matemática e Linguagens ; 8.
Refletir e compreender visões políticas de desconstrução da democracia que podem resultar em conflitos étnicos, religiosos, políticos, \textbf{econômicos} ligados à questão de sexualidade e gênero. \textbf{Espera} -se, assim, que a Sociologia, com suas \textbf{teorias} e conceitos, contribua, dentro do campo das Ciências Humanas, para uma \textbf{progressão} dos conhecimentos apresentados pelas \textbf{competências} e habilidades da Área do Conhecimento das Ciências Humanas, que vá do mais \textbf{simples} ao mais \textbf{complexo}, do mais concreto ao mais abstrato e \textbf{mobilize} o conjunto desses conceitos e \textbf{teorias} considerados \textbf{indispensáveis} para que as ( os ) estudantes possam pensar/refletir sobre a sociedade por meio da \textbf{ciência}.
Nesse novo arranjo curricular, destaca -se ainda o lugar que a Sociologia deve ocupar dentro da parte diversificada do Currículo Referência que é composta pelos Itinerários Formativos e as discussões sobre o “ Projeto de Vida ”.
Junto às discussões sobre o Projeto de Vida, os conhecimentos sobre a realidade social, diversidades e desigualdades, \textbf{mercado} de trabalho e outros temas abordados pelo componente curricular de Sociologia podem ajudar a complexificar os \textbf{debates} propostos nesse novo currículo do ensino \textbf{médio}.
Contribuindo para que esses \textbf{debates} não sejam uma \textbf{mera} discussão meritocrática e que encerrem em atitudes voluntaristas sem uma \textbf{devida} reflexão sobre o contexto social em que estão inseridos, mas que permitam evidenciar as dificuldades impostas pela estrutura social para que as ( os ) jovens alcancem seus objetivos nos mais diversos campos de sua vida.
Somente a partir dessas reflexões, os ( as ) estudantes podem adotar decisões assertivas em diferentes dimensões de suas vidas assim como \textbf{demandar} políticas para superação das diversas formas de desigualdades que \textbf{impactam} diretamente a efetivação de seus projetos de vida.
Já na construção dos Itinerários Formativos, a Sociologia com seus conhecimentos, \textbf{teorias} e \textbf{métodos} contribui para que as/os jovens pesquisem empiricamente a realidade social, política, cultural e \textbf{econômica} da sua comunidade.
O texto da BNCC estimula a construção de observatórios, núcleos de estudos, clubes de \textbf{debates}, núcleo de criação artísticas entre outras ações que permitam que a Sociologia amplie sua atuação dentro do ensino \textbf{médio} brasileiro promovendo não apenas a divulgação, mas também a produção \textbf{crítica} de conhecimento sobre a realidade social da comunidade onde está inserida a escola.

% matched lemmas: abordagem, análise, argumento, ator, aula, capacidade, cientista, ciência, competência, complexo, compreensão, conceitual, configurar, contemporâneo, contextualização, crítico, debate, demandar, depender, devido, disciplina, dizer, econômico, emergir, esperar, evento, exigir, explicitar, final, impactar, indispensável, interdisciplinaridade, maneira, mercado, mero, metodológico, mobilizar, moderno, médio, método, parâmetro, problema, progressão, reflexivo, simples, somar, subsídio, tanto, tarefa, tecnologia, teoria, teórico, totalidade, voz
\end{textsample}
